% !TEX program = xelatex
% =====================================================================
%  Chinese translation of ../slides/main.tex — keep frame structure in sync
%  Compile: ttx main   (xelatex + biber + xelatex x2)
% =====================================================================
% \documentclass[draft]{beamer}
\documentclass[]{beamer}
\usetheme{Berlin}
\usepackage{hyperref}
\usepackage{graphicx}   % figures
\usepackage{multicol}   % multi-column layouts
\usepackage{booktabs}   % tables
\usepackage{subfigure}  % sub-figures
\usepackage{xcolor}
\usepackage{color}
\usepackage{csquotes}
\usepackage{xltxtra}
\usepackage{xelatexemoji}
\usepackage[style=apa,sortcites=true,sorting=nyt,backend=biber]{biblatex}
\addbibresource{refs.bib}
\usepackage{tikz}       % DAG graphs
\usetikzlibrary{shapes,decorations,arrows,calc,arrows.meta,fit,positioning}

\setbeamertemplate{caption}{\raggedright\insertcaption\par}
\setbeamertemplate{footline}{}
\tikzset{
  -Latex,auto,node distance =1 cm and 1 cm,semithick,
  state/.style ={ellipse, draw, minimum width = 0.7 cm},
  block/.style ={rectangle, draw, minimum width = 0.7 cm},
  point/.style = {circle, draw, inner sep=0.04cm,fill,node contents={}},
  bidirected/.style={Latex-Latex},
  el/.style = {inner sep=2pt, align=left, sloped}
}

% ---------- Colours ----------
\definecolor{green}{RGB}{62, 138, 46}
\definecolor{blue}{RGB}{3, 105, 170}
\definecolor{red}{RGB}{194, 54, 55}

\usecolortheme{beaver}

\definecolor{primary}{RGB}{0,24,49}
\definecolor{secondary}{RGB}{0,48,96}
\definecolor{tertiary}{RGB}{94,48,0}
\definecolor{quaternary}{RGB}{94,48,0}
\setbeamercolor{titlelike}{parent=palette secondary,fg=white}
\setbeamercolor{frametitle}{bg=primary!60!white}
\setbeamercolor{frametitle right}{bg=primary!60!white}
\setbeamercolor*{palette primary}{bg=primary, fg = white}
\setbeamercolor*{palette secondary}{bg=secondary, fg = white}
\setbeamercolor*{palette tertiary}{bg=tertiary, fg = white}
\setbeamercolor*{palette quaternary}{bg=quaternary, fg = white}

\definecolor{nus_blue}{rgb}{0,0.24,0.49}
\definecolor{nus_yellow}{rgb}{0.94, 0.48, 0.00}
\definecolor{nus_yellow_highlight}{rgb}{0.94, 0.48, 0.00}

\setbeamercolor{section in toc}{fg=black,bg=white}
\setbeamercolor{alerted text}{fg=nus_blue!80!gray}
\setbeamercolor{navigation symbols}{fg=gray!60!black, bg=gray!10!white}
\setbeamercolor{navigation symbols dimmed}{fg=gray!50!black}
\setbeamercolor*{palette primary}{fg=black!60!black,bg=gray!15!white}
\setbeamercolor*{palette secondary}{fg=black!70!black,bg=gray!15!white}
\setbeamercolor*{palette tertiary}{bg=black!80!black,fg=gray!10!white}
\setbeamercolor*{palette quaternary}{fg=black,bg=gray!5!white}
\setbeamercolor{enumerate item}{fg=black,bg=white}
\setbeamercolor{itemize item}{fg=black,bg=white}
\setbeamercolor{item}{fg=black,bg=white}
\setbeamercolor*{sidebar}{fg=black,bg=gray!15!white}

\setbeamercolor*{palette sidebar primary}{fg=black!10!black}
\setbeamercolor*{palette sidebar secondary}{fg=white}
\setbeamercolor*{palette sidebar tertiary}{fg=black!50!black}
\setbeamercolor*{palette sidebar quaternary}{fg=gray!10!white}

\setbeamercolor{button}{bg=white,fg=black}
\setbeamercolor{author in head/foot}{bg=black}
\setbeamercolor{title in head/foot}{bg=black!20}
\setbeamercolor{titlelike}{bg=black!100!white}
\setbeamercolor{frametitle}{bg=black!100!white}
\setbeamercolor{frametitle right}{bg=black!100!white}
\setbeamercolor{empty}{bg=black}

% ---------- Headline (section navigation) & frame title ----------
\setbeamertemplate{headline}{%
  \begin{beamercolorbox}[wd=\paperwidth,ht=3.5ex,dp=3.5ex]{section in head/foot}
    \insertnavigation{\paperwidth}%
  \end{beamercolorbox}%
}

\setbeamertemplate{frametitle}{%
  \nointerlineskip%
  \begin{beamercolorbox}[wd=\paperwidth,leftskip=.3cm,rightskip=.3cm plus1fil,vmode]{frametitle}
    \vskip0.5ex%
    \usebeamerfont*{frametitle}\strut\insertframetitle\strut\par%
    \vskip0.05ex%
  \end{beamercolorbox}%
}

\setbeamercolor*{separation line}{}
\setbeamercolor*{fine separation line}{}
\setbeamercolor{bibliography entry author}{fg=black}
\setbeamercolor{bibliography entry title}{fg=black}
\setbeamercolor{bibliography entry location}{fg=black}
\setbeamercolor{bibliography entry note}{fg=black}
\setbeamertemplate{bibliography item}{}

% ---------- CJK support ----------
\usepackage{xeCJK}
\setCJKsansfont{PingFang SC}
\setCJKmainfont{Songti SC}

% ---------- Helper macros ----------
\providecommand{\tabfnm}[1]{\textsuperscript{#1}}  % apa7 marker used by \input tables
\newcommand{\hlt}[1]{\textcolor{red}{\textbf{#1}}}
\newcommand{\citegray}[1]{{\tiny\textcolor{gray}{\parencite{#1}}}}
\newcommand{\referencesize}{\tiny}
\renewcommand*{\bibfont}{\referencesize}

% Black divider slide. \sectionframe{Title} also starts a \section, so the
% headline navigation highlights correctly. Use \dividerframe{Title} for a
% divider without a new section (e.g. Acknowledgments, References).
\newcommand{\dividerframe}[1]{%
  {\setbeamercolor{background canvas}{bg=black}%
   \begin{frame}
     \begin{center}\Huge{\textcolor{white}{#1}}\end{center}
   \end{frame}}%
}
\newcommand{\sectionframe}[1]{\section{#1}\dividerframe{#1}}

% ---------- Title page ----------
\title[短标题]{在定量社会科学研究中负责任地使用人工智能}
\author[Rongxin Ouyang]{Presenter Name$^{1\#}$, Coauthor Name$^{1*}$}
\institute{\textcolor{gray}{
    $^{\#}$ 报告人。邮箱：presenter@example.com\\
    $^{*}$ 通讯作者。邮箱：corresponding@example.com\\
    ~\\
  $^{1}$单位，城市\\}}
\date{\small \textcolor{gray}{\the\year 年\the\month 月\the\day 日}}

\begin{document}

\frame{\titlepage}

% =====================================================================
\sectionframe{引言}
% =====================================================================

\begin{frame}
  \frametitle{引言：研究动机}
  \begin{itemize}
    \item 第一个要点，附灰色引用~\citegray{king2013how}
    \item 第二个要点，含\hlt{高亮}短语
      \begin{itemize}
        \item 子要点~\citegray{roberts2018censored}
        \item 中文文献来源~\citegray{weibo2026hotlist}
      \end{itemize}
  \end{itemize}
\end{frame}

% 左侧文字，右侧图片
\begin{frame}
  \frametitle{引言：文字与图片}
  \begin{columns}[T,onlytextwidth]
    \column{0.6\textwidth}
      \begin{itemize}
        \item 已有研究发现 X \parencite{king2013how}
        \item \textcolor{gray}{此前已展示的要点，灰显以逐步呈现}
      \end{itemize}
    \column{0.4\textwidth}
      \begin{figure}
        \centering
        \includegraphics[width=1\linewidth]{example-image-a} % replace with fig/xxx.png
        \caption{\tiny \textcite{king2013how}}
      \end{figure}
  \end{columns}
\end{frame}

\begin{frame}
  \frametitle{引言：研究思路}
  \begin{itemize}
    \item 我们研究\ldots
    \item 我们测量\ldots
    \item 我们识别\ldots
    \item 我们验证\ldots
  \end{itemize}
\end{frame}

% =====================================================================
\sectionframe{文献综述}
% =====================================================================

\begin{frame}
  \frametitle{文献 1：主题}
  \begin{columns}[T,onlytextwidth]
    \begin{column}{0.67\textwidth}
      \footnotesize\let\normalsize\footnotesize
      \begin{enumerate}
        \item \textbf{阶段一。}描述 \normalsize ~\citegray{king2013how}
        \item \textbf{阶段二。}描述 \normalsize ~\citegray{roberts2018censored}
      \end{enumerate}
    \end{column}
    \begin{column}{0.33\textwidth}
      \begin{figure}
        \centering
        \includegraphics[width=0.92\linewidth]{example-image-b}
        \caption{\tiny \textcite{roberts2018censored}}
      \end{figure}
    \end{column}
  \end{columns}
\end{frame}

\begin{frame}
  \frametitle{文献 1：研究假设}
  \small
  \begin{itemize}
    \item \textbf{H1. 名称。}假设陈述
    \item \textbf{H2. 名称。}假设陈述
  \end{itemize}
  ~\\
  \begin{center}
    \begin{tikzpicture}
      \node[block] (a) at (0,1.6) {\small 机制 A \tiny (H1)};
      \node[block] (b) at (0,0) {\small 机制 B \tiny (H2)};
      \node[state, minimum width=1.6cm] (y) at (7.2,0.8) {\small $Y$};
      \path (a) edge node[el,above]{\tiny 渠道 a} (y);
      \path (b) edge node[el,below]{\tiny 渠道 b} (y);
    \end{tikzpicture}
  \end{center}
\end{frame}

% =====================================================================
\sectionframe{数据}
% =====================================================================

\begin{frame}[shrink]
  \frametitle{数据：来源}
  \begin{enumerate}
    \item \textbf{数据集 A}：描述 \tiny （覆盖范围、频率） \normalsize
      \begin{itemize}
        \item N 个观测值；关键变量
      \end{itemize}
    \item \textbf{数据集 B}：描述
  \end{enumerate}
\end{frame}

% =====================================================================
\sectionframe{研究方法}
% =====================================================================

\begin{frame}
  \frametitle{方法 1：模型}
  \begin{itemize}
    \item 分析单位：\textbf{单位}
    \item 估计方法~\citegray{king2013how}
    \small
    \begin{itemize}
      \item $x_i$：核心解释变量
      \item $z_{it}$：控制变量
    \end{itemize}
  \end{itemize}
  \begin{equation}
    y_{it} = \alpha + \beta_1 \textcolor{red}{x_i} + \beta_2 z_{it} + \varepsilon_{it}
  \end{equation}
\end{frame}

% =====================================================================
\sectionframe{研究结果}
% =====================================================================

\begin{frame}
  \frametitle{结果：概览}
  \begin{enumerate}
    \item \textbf{发现 1}：\ldots
    \item \textbf{发现 2}：\ldots
    \item \textbf{发现 3}：\ldots
  \end{enumerate}
\end{frame}

% 回归表（可将 tabular 替换为 \input{../output/tab/xxx.tex}）
\begin{frame}[shrink]
  \frametitle{结果 1：主要估计}
  \begin{table}
    \resizebox{0.6\linewidth}{!}{%
      \begin{tabular}{lccc}
        \toprule
                & M1      & M2      & M3      \\
        \midrule
        $x_i$   & 0.040\tabfnm{***} & 0.038\tabfnm{***} & 0.034\tabfnm{***} \\
                & (0.003) & (0.003) & (0.003) \\
        控制变量 & 否     & 是     & 是     \\
        \bottomrule
      \end{tabular}
    }
  \end{table}
  \tiny \textcolor{gray}{注：$^{*} p < .05$, $^{**} p < .01$, $^{***} p < .001$.}
\end{frame}

% 左侧要点，右侧结果图
\begin{frame}
  \frametitle{结果 1：解读}
  \begin{columns}
    \column{0.4\textwidth}
    \footnotesize\let\normalsize\footnotesize
    \begin{itemize}
      \item 系数为正 \tiny \textcolor{gray}{（$.034$，$p < .001$）} \normalsize ，约为\hlt{每个标准差 3.5\%}
      \item 第二个要点
    \end{itemize}
    \column{0.6\textwidth}
    \centering
    \includegraphics[width=0.85\linewidth]{example-image-c}\par
    \tiny \textcolor{gray}{灰色图注。}
  \end{columns}
\end{frame}

% =====================================================================
\sectionframe{讨论}
% =====================================================================

\begin{frame}
  \frametitle{讨论}
  \begin{itemize}
    \item 启示 1
    \item 启示 2
  \end{itemize}
\end{frame}

\begin{frame}
  \frametitle{局限性}
  \begin{enumerate}
    \item \textbf{测量。}\ldots
    \item \textbf{识别。}\ldots
    \item \textbf{外部有效性。}\ldots
  \end{enumerate}
\end{frame}

% =====================================================================
\sectionframe{补充材料}
% =====================================================================

\begin{frame}
  \frametitle{补充材料：两个子图}
  \begin{figure}
    \subfigure[子图 A]{
      \begin{minipage}[l]{0.45\linewidth}
        \includegraphics[width=1\linewidth]{example-image-a}
      \end{minipage}
    }
    \subfigure[子图 B]{
      \begin{minipage}[l]{0.45\linewidth}
        \includegraphics[width=1\linewidth]{example-image-b}
      \end{minipage}
    }
  \end{figure}
\end{frame}

{
  \setbeamercolor{background canvas}{bg=black}
  \begin{frame}
    \begin{center}
      \Huge{\textcolor{white}{致谢}}
    \end{center}
    \vspace{1cm}
    \normalsize{\textcolor{white}{感谢\ldots。文责自负。}}
  \end{frame}
}

\dividerframe{参考文献}

\begin{frame}[allowframebreaks]
  \frametitle{参考文献}
  \begingroup
  \referencesize\color{black}
  \printbibliography[notkeyword=chinese,heading=none]
  \endgroup
\end{frame}

% refs.bib 中标注 `keywords = {chinese}` 的条目列于此页
\begin{frame}[allowframebreaks]
  \frametitle{参考文献（中文）}
  \begingroup
  \referencesize\color{black}
  \printbibliography[keyword=chinese,heading=none]
  \endgroup
\end{frame}

\end{document}
